\documentclass[10pt,a4paper]{amsart}
\usepackage[margin=19mm]{geometry}
\usepackage{amsmath,amssymb,mathtools}
\usepackage[scr=rsfs,cal=euler]{mathalfa}
\usepackage{xfrac}
\usepackage[unicode,hidelinks,bookmarks=false]{hyperref}
\newcommand{\ie}{i.e.\ }
\newcommand{\cf}{cf.\ }
\newcommand{\resp}{respectively\ }
\newcommand{\Subsection}{\subsection}
\providecommand{\nobreakdash}{\nobreak\hbox{-}}
\setlength{\emergencystretch}{2em}
\multlinegap0pt
\usepackage{bm}
\usepackage{bbm}
\usepackage{dsfont}
\usepackage{xspace}
\usepackage{cancel}
\usepackage{mathtools}
\usepackage{amsthm}
\makeatletter
\@ifundefined{prop}{\newtheorem{prop}{Prop}}{}
\@ifundefined{lemma}{\newtheorem{lemma}[prop]{Lemma}}{}
\makeatother
\title[Compact Moduli Spaces of Marked Cubic Plane Curves]{Compact Moduli Spaces of Marked Cubic Plane Curves}
\author{Aaron Goodwin}
\date{Source: arXiv:2501.14927v2 (CC BY 4.0)}
\begin{document}
\maketitle

\noindent\emph{Source and licence.} Compact Moduli Spaces of Marked Cubic Plane Curves. Version arXiv:2501.14927v2 (CC BY 4.0), distributed under the Creative Commons Attribution 4.0 International licence, as recorded in the arXiv licence metadata for this exact version and shown on the abstract page https://arxiv.org/abs/2501.14927v2. Full rights notice: licenses/arxiv-2501.14927-CC-BY-4.0.txt.\par\smallskip

\noindent\textit{Editorial scope.} Counted unit: the Lemma whose source label is \texttt{Lemma:wallC(A\_3)} in the article (source0.tex lines 1202--1204, source label \texttt{Lemma:wallC(A\_3)}), delivered together with the article's own complete original proof of it (source0.tex lines 1206--1387, one \texttt{proof} environment of 182 lines). No mathematics is written, repaired or re-derived: statement and proof are the article's own text line for line. Required context retained uncounted: lemma:Wall\_3distinctlines (lines 818--823). Every definition, display and label of those lines is retained unchanged. Omitted from this package and not redistributed: 1--817, 824--1201, 1205--1205, 1388--1784. Nothing inside a retained excerpt is omitted or altered: every retained body is delivered exactly as the article's lines are written, so no omission receipt of any kind accompanies this package. Citations: the retained excerpts carry no citation command at all, so no bibliography entry of the article is transcribed and no reference is redistributed. Reference adaptation: none was needed and none was made; every reference inside the delivered unit resolves inside it, so no unresolved reference or citation is produced. Printed slips kept literal: no printed word, symbol, hypothesis or number of the counted statement or of its proof was altered. Layout and class: the article's own class is \texttt{amsart} with packages including inputenc, url, hyperref, geometry, amsmath, amsthm, amssymb, graphicx, multicol, array, mathtools, mathrsfs, tikz-cd, yfonts; the wrapper is the lane's pinned \texttt{amsart} editorial wrapper at 10pt on A4, which restates the theorem-like environments the retained text needs and adds the article's own macro definitions that the retained text needs (none), with extra wrapper packages (bm, bbm, dsfont, xspace, cancel, mathtools). The delivered numbering is the unit's own. Assets: no retained span contains a figure environment, a graphics-inclusion command, a PDF-page inclusion, a table or a TikZ picture; no image file is copied into this package, the package declares no asset dependency and no image file is redistributed with it. Licence: version arXiv:2501.14927v2 of this article is distributed under the Creative Commons Attribution 4.0 International licence, as recorded in the arXiv licence metadata for this exact version and shown on the abstract page https://arxiv.org/abs/2501.14927v2. A full rights notice is delivered at licenses/arxiv-2501.14927-CC-BY-4.0.txt. Characters: every retained excerpt is pure ASCII, so the delivered engine, XeLaTeX with its default Latin Modern text font, reports no missing character.
\begin{lemma}[Case of $C(3A_1)$]
\label{lemma:Wall_3distinctlines}
    Let $\mathbf{x}:=(C, p_1, \dots p_n)$ have positive dimensional stabilizer, where the curve $C$ is three lines in general position. Let $\eta$ and $\Lambda$ be as above. Define $I$ to be the subset of $\{1, \dots , n\}$ such that the points $p_i$ coincide at $\eta \in C$ for $i \in I$ and the remaining points lie on the opposite line $\Lambda \subset C$. Then $\mathbf{x}$ is $L$ strictly semi-stable if and only if $\sum_{i \in I}w_i = \frac{1}{2} \sum_{j \notin I} w_j$, the total weight at each of the other nodes is at most $\frac{1}{3} \sum_1^n w_i$, and the weight at any smooth point of $\Lambda$  is at most $\frac{1}{3}\sum_1^n w_i + \gamma $.

    
\end{lemma}

\begin{lemma}[Case of $C(A_3)$] \label{Lemma:wallC(A_3)}
    Let $\mathbf{x} := (C, p_1, \dots , p_n)$ have positive dimensional stabilizer, where $C$ is the union of a conic $A$ and a tangent line $B$. Let $I$ be the subset of $\{1, \dots ,n\}$ such that $p_i$ coincides with the tacnode for $i \in I$ and let $J$ be the subset such that the $p_j$ coincide at a non-singular point of $B$ for $j \in J$. The remaining marked points coincide with the point on the conic $A$ whose tangent line intersects the $p_j$s. Then $\mathbf{x}$ is $L$ strictly semi-stable if and only if $\sum_{k \notin I \cup J} w_k = \sum_{i \in I} w_i + \gamma$ and $\sum_{i \in I} w_i - \gamma \leq \sum_{j \in J} w_j \leq \sum_{i \in I} w_i + 2 \gamma$.
\end{lemma}

\begin{proof}
First we remark that if all of the marked points lie on the line $B$ then $\mathbf{x}$ is unstable for all line bundles $L$. This can be seen by letting $g \in SL(3)$ take the tacnode to $(1:0:0)$ and the line $B$ to $\{z=0\}$. Then $\mu^L(g \cdot \mathbf{x}, \lambda_1) = \sum_1^n w_i$, which is positive.

Now assume that the marked points lie either at the tacnode or on the intersection $C \cap \Lambda$ where $\Lambda$ is a line \textit{other than $B$} which is tangent to the conic. We proceed as in Lemma \ref{lemma:Wall_3distinctlines} by solving $\max \{ \mu^L(g \cdot \mathbf{x}, \lambda_r) \ | \ g \in SL(3), \ r \in [\frac{-1}{2}, 1] \} = 0$ for $L$. There are  six subsets of $SL(3)$ which maximize $\Xi_{max}$. 

If $g\in SL(3)$ maximizes $\Xi_{max}(g \cdot \mathbf{x})$ then it must take the tacnode or one of other the marked points to $(1:0:0)$. Let $p_t$ be the tacnode, $p_a$ the marked point on the conic, $p_b$ the marked point on the linear component of $C$, and $\Lambda$ the line through $p_a$ and $p_b$. If $g_1$ takes $p_t$ to $(1:0:0)$ and $B$ to $\{z=0\}$ then $\Xi_{max}( g \cdot C) = \{xz^2, y^2z\}$. If $g_2$ takes $p_b$ to $(1:0:0)$ and $B$ to $\{z=0\}$ then $\Xi_{max} ( g_2 \cdot \mathbf{x}) = \{x^2z,  xy^2 \}$. If $g_3$ takes $p_b$ to $(1:0:0)$ and $\Lambda$ to $\{z=0\}$ then $\Xi_{max} ( g_3 \cdot \mathbf{x}) = \{x^2y\}$. If $g_4$ takes  $p_a$ to $(1:0:0)$ and $\Lambda$ to $\{z=0\}$ then $\Xi_{max} ( g_4 \cdot \mathbf{x}) = \{ x^2z, xy^2\}$. If $g_5$ takes $p_t$ to $(1:0:0)$ and the line $\overline{p_a p_t}$ to $\{z=0\}$ then $\Xi_{max} ( g_5 \cdot \mathbf{x}) = \{ xy^2\}$. Finally, if $g_6$ takes $p_a$ to $(1:0:0)$ and the line $\overline{p_a p_t}$ to $\{z=0\}$ then $\Xi_{max} ( g_6 \cdot \mathbf{x}) = \{x^2y\}$.

We abbreviate $\sum_{p_i = p_t} w_i$ as $w_t$, $\sum_{p_j = p_b} w_j$ as $w_b$, and $\sum_{p_k = p_a} w_k$ as $w_a$. Then the maximal values of $\mu^L( g \cdot \mathbf{x}, \lambda_r)$ are

\begin{itemize}
    \item $\mu^L(g_1 \cdot \mathbf{x}, \lambda_r) = w_t + rw_b +(-1-r)w_a - \gamma \max\{ -2-r, -1+r \}$, which is maximized in $r$ by
     \begin{align*}
        F_1 := \mu^L(g_1 \cdot \mathbf{x}, \lambda_1) = p_t -2p_a +p_b
        && \text{and} &&
        F_2 := \mu^L(g_1 \cdot \mathbf{x}, \lambda_0) = p_t - p_a + \gamma \end{align*}
        \begin{align*} \text{and} \ \ F_3 := \mu^L(g_1 \cdot \mathbf{x}, \lambda_{\frac{-1}{2}}) = p_t - \frac{p_a}{2} - \frac{p_b}{2}.   \end{align*}
    
    
    \item $\mu^L(g_2 \cdot \mathbf{x}, \lambda_r) = w_b + rw_t +(-1-r) w_a - \gamma \max\{1-r, 1+2r \}$, which is maximized in $r$ by
      \begin{align*}
        F_4 := \mu^L(g_2 \cdot \mathbf{x},  \lambda_1) = p_b -2 p_a + p_t + 3 \gamma 
        && \text{and} &&
        F_5 := \mu^L(g_2 \cdot \mathbf{x}, \lambda_{\frac{-1}{2}}) = p_b - \frac{1}{2}p_a - \frac{1}{2}p_t
    \end{align*}
    
    \item $\mu^L(g_3 \cdot \mathbf{x}, \lambda_r) = w_b + rw_a + (-1-r)w_t - \gamma (2+r)$, which is maximized in $r$ by
     \begin{align*}
        F_6 := \mu^L(g_3 \cdot \mathbf{x}, \lambda_1) = p_b -2 p_a + p_t + 3 \gamma
        && \text{and} &&
        F_7 := \mu^L(g_3 \cdot \mathbf{x}, \lambda_{\frac{-1}{2}}) = p_b - \frac{p_a}{2} - \frac{p_t}{2}.
    \end{align*}
    
    \item $\mu^L(g_4 \cdot \mathbf{x}, \lambda_r) = w_a + r w_b + (-1-r) w_t - \gamma \max \{1-r, 1+2r\}$, which is maximized in $r$ by
    \begin{align*}
        F_8 := \mu^L(g_1 \cdot \mathbf{x}, \lambda_1) = p_a + p_b - 2 p_t - 3 \gamma
        && \text{and} &&
        F_9 := \mu^L(g_1 \cdot \mathbf{x},  \lambda_0) = p_a - p_t - \gamma
        \end{align*}
        \begin{align*} \text{and} \ \
        F_{10} := \mu^L(g_1 \cdot \mathbf{x}, \lambda_{\frac{-1}{2}}) = p_a - \frac{p_b}{2} - \frac{p_t}{2} - \frac{3\gamma}{2}.
    \end{align*}
    
    
    \item $\mu^L(g_5 \cdot \mathbf{x}, \lambda_r) = w_t + rw_a + (-1-r)w_b - \gamma(1+2r)$, which is maximized in $r$ by
    \begin{align*}
        F_{11} := \mu^L(g_5 \cdot \mathbf{x}, \lambda_1) = w_t + w_a - 2w_b -3\gamma
        && \text{and} &&
        F_{12} := \mu^L(g_5 \cdot \mathbf{x}, \lambda_{\frac{-1}{2}}) = w_t - \frac{-1}{2}w_a - \frac{1}{2}w_b
    \end{align*}
    
    \item $\mu^L(g_6 \cdot \mathbf{x}, \lambda_r) = w_a + rw_t + (-1-r)w_b - \gamma(2+r)$, which is maximized in $r$ by
    \begin{align*}
        F_{13} := \mu^L(g_6 \cdot \mathbf{x}, \lambda_1) = w_a -2w_b + w_t - 3\gamma
        && \text{and} &&
        F_{14} := \mu^L(g_6 \cdot \mathbf{x}, \lambda_{\frac{-1}{2}}) = w_a - \frac{1}{2}w_b - \frac{1}{2}w_t - \frac{3}{2}\gamma.
    \end{align*}
\end{itemize}
We solve $\max \{F_i\} =0$ for $(\gamma, w_1, \dots , w_n)$. Since $F_2 = -F_9,$ they must both be zero. This gives the wall $p_a = p_t + \gamma.$ Evaluating the remaining inequalities at this wall gives the condition $w_t - \gamma \leq w_b \leq w_t +2 \gamma.$




% \vspace{2cm}

% We analyze four cases [MISSING CASES? - $g_5 and g_6$ do not alter result]: If $g$ takes $B$ to $\{z=0\}$ then $\mu$ is maximized when $g$ takes $p_t$ or $p_b$ to $(0:0:1)$. In Case 1 $g$ takes $p_t$ to $(0:0:1)$ and in Case 2 $g$ takes $p_b$ to $(0:0:1)$. If $g$ does not take $B$ to $\{z=0\}$ then $\mu$ is maximized by $g \in SL(3)$ which take the line through $p_a$ and $p_b$ to $\{z=0\}$ and either $p_a$ or $p_b$ to $(1:0:0)$. In Case 3 $g$ takes $p_b$ to $(1:0:0)$ and in Case 4 $g$ takes $p_a$ to $(1:0:0).$

% \vspace{0.5cm}

% Case 1: Assume $g$ takes $B$ to $\{z=0\}$ and $p_t$ to $(1:0:0)$. Then the maximal support of $g \cdot C$ is $\{xz^2, y^2z\}$ and

% $$\mu^L (g \cdot (C, p_1, \dots , p_n), \lambda_r) = p_t + rp_b + (-1-r) p_a - \gamma \max \begin{cases}
%     -1-2r \\
%     -1 + r
% \end{cases}$$
% $$ = p_t + rp_b + (-1-r) p_a - \gamma \begin{cases}
%     -1-2r, \ \ \  r \leq 0\\
%     -1 + r, \ \ \ r \geq 0 
% \end{cases}$$
% $$ = p_t - p_a + \gamma + \begin{cases}
%     r(p_b - p_a + 2 \gamma), \ \ \ r \leq 0 \\
%     r(p_b - p_a - \gamma), \ \ \ r \geq 0
% \end{cases}.$$

% This is maximized by $r=1$ if $p_b \geq p_a + \gamma$; by $r=0$ if $ p_a - 2 \gamma \leq p_b \leq p_a + \gamma;$ and by $r= \frac{-1}{2}$ if $p_b \leq p_a - 2 \gamma$.

% $$\mu^L (g \cdot (C, p_1, \dots , p_n), \lambda_1) = p_t -2p_a +p_b. $$





% $$\mu^L (g \cdot (C, p_1, \dots , p_n), \lambda_0) = p_t - p_a + \gamma.$$






% $$\mu^L (g \cdot (C, p_1, \dots , p_n), \lambda_{\frac{-1}{2}}) = p_t - \frac{p_a}{2} - \frac{p_b}{2}. $$


% \vspace{0.5cm}

% Case 2: Assume $g$ takes $B$ to $\{z=0\}$ and $p_b$ to $(1:0:0)$. Then the maximal support of $g \cdot C$ is $\{ xz^2\}$. Then 

% $$ \mu^L (g \cdot (C, p_1, \dots , p_n), \lambda_r) = p_b + rp_t + (-1-r)p_a - \gamma (-1-2r ) $$
% $$ = p_b - p_a + \gamma + r(p_t - p_a + 2 \gamma). $$

% This is maximized by $r = 1$ or $r= \frac{-1}{2}$ depending on whether $p_t + 2\gamma \leq p_a$. 

% $$ \mu^L (g \cdot (C, p_1, \dots , p_n), \lambda_1) = p_b -2 p_a + p_t + 3 \gamma .$$




% $$\mu^L (g \cdot (C, p_1, \dots , p_n), \lambda_{\frac{-1}{2}}) = p_b - \frac{p_a}{2} - \frac{p_t}{2} .$$


% \vspace{0.5cm}

% Case 3: Assume $g$ takes the line through $p_a$ and $p_b$ to $\{z=0\}$ and $p_b$ to $(1:0:0).$ Then the maximal support of $g \cdot C$ is $\{x^2z\}$ and
% $$\mu^L( g \cdot (p_1 , \dots , p_n), \lambda_r) = p_b + rp_a + (-1-r)p_t - \gamma (2+r) $$
% $$ = p_b - p_t - 2 \gamma + r(p_a - p_t - \gamma ).$$

% This is maximized by $r = 1$ or $r= \frac{-1}{2}$, depending on whether $p_a  \leq p_t + \gamma$.

% $$\mu^L( g \cdot (p_1 , \dots , p_n), \lambda_1) = p_b + p_a - 2p_t - 3\gamma  .$$




% $$\mu^L( g \cdot (p_1 , \dots , p_n), \lambda_{\frac{-1}{2}}) = p_b - \frac{p_a}{2} - \frac{p_t}{2} - \frac{3 \gamma}{2}. $$

% \vspace{0.5cm}

% Case 4: Assume $g$ takes the line through $p_a$ and $p_b$ to $\{z=0\}$ and $p_a$ to $(1:0:0)$. Then the maximal support of $g \cdot C$ is $\{x^2z,  xy^2 \}$ and

% $$ \mu^L( g \cdot (p_1 , \dots , p_n), \lambda_r) = p_a  + rp_b + (-1-r)p_t - \gamma \max \begin{cases}
%     1-r \\
%     1+2r
% \end{cases}$$
% $$ = p_a  + rp_b + (-1-r)p_t - \gamma \begin{cases}
%     1-r, \ \ \ r \leq 0 \\
%     1+2r, \ \ \ r \geq 0
% \end{cases}$$

% $$ = p_a - p_t - \gamma + \begin{cases}
%     r( p_b - p_t + \gamma), \ \ \ r \leq 0 \\
%     r(p_b - p_t - 2 \gamma), \ \ \ r \geq 0
% \end{cases}.$$

% This is either maximized by $r= 1$, $r= 0$, or $r = \frac{-1}{2}.$

% $$  \mu^L( g \cdot (p_1 , \dots , p_n), \lambda_1) = p_a + p_b - 2 p_t - 3 \gamma.$$



% $$  \mu^L( g \cdot (p_1 , \dots , p_n), \lambda_0) = p_a - p_t - \gamma.$$




% $$  \mu^L( g \cdot (p_1 , \dots , p_n), \lambda_{\frac{-1}{2}}) = p_a - \frac{p_b}{2} - \frac{p_t}{2} - \frac{3\gamma}{2}. $$

% \vspace{0.5cm}
%  The  strict semi-stability of $(C, p_1, \dots , p_n)$ then implies that the maximum of the following maximal values of $\mu^L(g \cdot (C, p_1, \dots , p_n), \lambda_r)$ found above is equal to $0$:

%  \begin{enumerate}
%      \item $p_t - 2 p_a + p_b$
%      \item $p_t - p_a + \gamma$
%      \item $p_t - \frac{p_a}{2} - \frac{p_b}{2}$
%      \item $p_b - 2p_a + p_t + 3 \gamma$
%      \item $p_b - \frac{p_a}{2} - \frac{p_t}{2}$
%      \item $p_b + p_a - 2p_t -3 \gamma$
%      \item $p_b - \frac{p_a}{2} - \frac{p_t}{2} - \frac{3 \gamma}{2}$
%      \item $p_a - p_t - \gamma$
%      \item $p_a - \frac{p_b}{2} - \frac{p_t}{2} - \frac{3 \gamma}{2}.$
%  \end{enumerate}

%  Term 4 is less than term 3 and term 5 is less than term 7, so terms 4 and 5 are not among the maximal values of $\mu$. Term 8 is $-1$ times term 2, so the fact that they are both $\leq 0$ implies that $p_a = p_t + \gamma.$ When evaluated at $p_a = p_t + \gamma$ the remaining terms become either $-p_t - 2 \gamma + p_b$ or $ \frac{p_t}{2} - \frac{\gamma}{2} - \frac{p_b}{2}$. It follows that $p_t - \gamma \leq p_b \leq p_t - 2 \gamma.$
    
\end{proof}

\end{document}
